\section*{HISTORICAL NOTE}
\phantomsection
\addcontentsline{toc}{section}{Historical Note}

(Numbers in brackets refer to the bibliography at the end of this Note.)

As we remarked in the Historical Note to Chapter II, the notion of a metric space was introduced in 1906 by M. Fréchet, and developed some years later by F. Hausdorff in his ``Mengenlehre''. It acquired great importance after 1920, partly as a consequence of the fundamental work of S. Banach and his school on normed spaces and their applications to functional analysis, and also because of the significance of the notion of absolute value in the theory of numbers and algebraic geometry (where in particular the process of completion with respect to an absolute value has proved to be a powerful instrument).

The decade 1920-1930 saw a whole series of investigations into the properties of metric spaces. These studies, undertaken by the Moscow school, aimed especially at getting necessary and sufficient conditions for the metrizability of a given topology. This movement of ideas brought out the significance of the notion of a normal space, which had been defined by Tietze in 1923 but whose important role was recognized only as a consequence of the work of Urysohn {[}6{]} on the extension of continuous real-valued functions. Except for the trivial case of functions of a real variable, the problem of extending to the whole space a continuous real-valued function defined on a closed set was first considered (for the case of the plane) by H. Lebesgue {[}3{]}; before Urysohn's definitive result, it was solved for metric spaces by H. Tietze {[}4{]}. The generalization of this problem to include functions with values in an arbitrary topological space has acquired considerable importance in algebraic topology in the last few years. Recent work has also shown that in questions of this sort the notion of a normal space is not easily handled, since it allows too many possibilities of ``pathology''; it often has to be replaced by the more restrictive concept of paracompactness, introduced in 1944 by J. Dieudonné {[}9{]}. In this theory, the most remarkable result is the theorem of A. H. Stone {[}10{]} which states that every metrizable space is paracompact (*).

We have already noted (Historical Note to Chapter IV) the important work at the end of the nineteenth century and the beginning of the twentieth (E. Borel, Baire, Lebesgue, Osgood, W. H. Young) on the classification of point-sets in \(R^{n}\) , and on the classification and characterization of real-valued functions obtained from continuous functions by iterating the process of passing to the limit (with respect to sequences of functions). It was quickly realized that metric spaces, whose development after 1910 is largely the work of the Russian and Polish schools, formed a natural domain for investigations of this nature. Among other things, the study of metric spaces has shown clearly the fundamental role played in modern analysis by the notion of a meagre set and by the theorem on the countable intersection of dense open sets in a complete metric space ( \(§\ 5,\ no.\ 3,\ Theorem\ 1\) ), which was first proved (independently) by Osgood {[}1{]} for the real line and by Baire {[}2{]} for the spaces \(R^{n}\) .

On the other hand, Souslin in 1917 {[}5{]}, correcting an error of Lebesgue's, showed that the continuous image of a Borel set is not necessarily a Borel set; this led him to the definition and study of a larger class of sets, since called ``analytic sets'' or ``Souslin sets''. After Souslin's premature death this work was carried forward particularly by N. Lusin (whose ideas had inspired Souslin's work) and the Polish mathematicians (see {[}7{]} and {[}8{]}). The importance of these sets nowadays lies in their applications to the theory of integration (where, thanks to their special properties, they allow constructions which would be impossible for arbitrary measurable sets) and to modern potential theory, in which the fundamental theorem on the capacitability of Souslin sets, proved recently by G. Choquet {[}11{]}, has already shown itself to be rich in diverse applications.

